% BEGIN GENERATED ARTIFACT: def:literal-propositional-logic
\begin{tcolorbox}[colback=propbox, colframe=propborder, arc=2pt,
  left=6pt, right=6pt, top=4pt, bottom=4pt,
  title={\small\textbf{Definition (Literal)}},
  fonttitle=\small\bfseries]
\begin{definition}[Literal]
\label{def:literal-propositional-logic}
A \emph{literal} is either a propositional variable \(P\) or the negation
\(\neg P\) of a propositional variable \(P\).
\end{definition}
\end{tcolorbox}

\begin{remark*}[Standard quantified statement]
For a formula \(\ell\in\WFF\),
\[
\ell\text{ is a literal}
\Longleftrightarrow
(\exists P\in\Prop)(\ell=P\text{ or }\ell=\neg P).
\]
\end{remark*}

\begin{remark*}[Definition predicate reading]
\[
\operatorname{Literal}(\ell)
\]
means that \(\ell\) is an atomic propositional variable or the negation of one
atomic propositional variable.
\end{remark*}

\begin{remark*}[Interpretation]
Literals are the indivisible signed atoms used in normal forms. They remember
which propositional variable is being tested and whether it appears positively
or negatively.
\end{remark*}

\begin{remark*}[Examples]
The formulas \(P\), \(Q\), \(\neg P\), and \(\neg Q\) are literals.
\end{remark*}

\begin{remark*}[Non-Examples]
The formulas \(P\land Q\), \(P\lor Q\), and \(\neg(P\land Q)\) are not
literals. The last formula is negated, but it negates a compound formula rather
than a single propositional variable.
\end{remark*}

\begin{dependencies}
\begin{itemize}
  \item \hyperref[def:propositional-variable]{Propositional Variable}
  \item \hyperref[def:logical-connective]{Logical Connective}
\end{itemize}
\end{dependencies}
% END GENERATED ARTIFACT: def:literal-propositional-logic

% BEGIN GENERATED ARTIFACT: def:positive-literal-propositional-logic
\begin{tcolorbox}[colback=propbox, colframe=propborder, arc=2pt,
  left=6pt, right=6pt, top=4pt, bottom=4pt,
  title={\small\textbf{Definition (Positive Literal)}},
  fonttitle=\small\bfseries]
\begin{definition}[Positive Literal]
\label{def:positive-literal-propositional-logic}
A \emph{positive literal} is a literal of the form \(P\), where \(P\) is a
propositional variable.
\end{definition}
\end{tcolorbox}

\begin{remark*}[Standard quantified statement]
For a formula \(\ell\in\WFF\),
\[
\ell\text{ is a positive literal}
\Longleftrightarrow
(\exists P\in\Prop)(\ell=P).
\]
\end{remark*}

\begin{remark*}[Definition predicate reading]
\[
\operatorname{PositiveLiteral}(\ell)
\]
means that \(\ell\) is a literal whose sign is positive.
\end{remark*}

\begin{remark*}[Interpretation]
A positive literal asserts the propositional variable directly. It carries no
leading negation.
\end{remark*}

\begin{remark*}[Examples]
The formulas \(P\) and \(Q\) are positive literals.
\end{remark*}

\begin{dependencies}
\begin{itemize}
  \item \hyperref[def:literal-propositional-logic]{Literal}
\end{itemize}
\end{dependencies}
% END GENERATED ARTIFACT: def:positive-literal-propositional-logic

% BEGIN GENERATED ARTIFACT: def:negative-literal-propositional-logic
\begin{tcolorbox}[colback=propbox, colframe=propborder, arc=2pt,
  left=6pt, right=6pt, top=4pt, bottom=4pt,
  title={\small\textbf{Definition (Negative Literal)}},
  fonttitle=\small\bfseries]
\begin{definition}[Negative Literal]
\label{def:negative-literal-propositional-logic}
A \emph{negative literal} is a literal of the form \(\neg P\), where \(P\) is a
propositional variable.
\end{definition}
\end{tcolorbox}

\begin{remark*}[Standard quantified statement]
For a formula \(\ell\in\WFF\),
\[
\ell\text{ is a negative literal}
\Longleftrightarrow
(\exists P\in\Prop)(\ell=\neg P).
\]
\end{remark*}

\begin{remark*}[Definition predicate reading]
\[
\operatorname{NegativeLiteral}(\ell)
\]
means that \(\ell\) is a literal whose sign is negative.
\end{remark*}

\begin{remark*}[Interpretation]
A negative literal denies exactly one propositional variable. In normal forms,
negation is pushed down until it appears only at this literal level.
\end{remark*}

\begin{remark*}[Examples]
The formulas \(\neg P\) and \(\neg Q\) are negative literals.
\end{remark*}

\begin{dependencies}
\begin{itemize}
  \item \hyperref[def:literal-propositional-logic]{Literal}
\end{itemize}
\end{dependencies}
% END GENERATED ARTIFACT: def:negative-literal-propositional-logic

% BEGIN GENERATED ARTIFACT: def:complementary-literals-propositional-logic
\begin{tcolorbox}[colback=propbox, colframe=propborder, arc=2pt,
  left=6pt, right=6pt, top=4pt, bottom=4pt,
  title={\small\textbf{Definition (Complementary Literals)}},
  fonttitle=\small\bfseries]
\begin{definition}[Complementary Literals]
\label{def:complementary-literals-propositional-logic}
Two literals are \emph{complementary} if they are \(P\) and \(\neg P\) for the
same propositional variable \(P\).
\end{definition}
\end{tcolorbox}

\begin{remark*}[Standard quantified statement]
For literals \(\ell_1,\ell_2\),
\[
\ell_1\text{ and }\ell_2\text{ are complementary}
\Longleftrightarrow
(\exists P\in\Prop)
\bigl((\ell_1=P\land \ell_2=\neg P)
\lor
(\ell_1=\neg P\land \ell_2=P)\bigr).
\]
\end{remark*}

\begin{remark*}[Definition predicate reading]
\[
\operatorname{ComplementaryLiterals}(\ell_1,\ell_2)
\]
means that the two literals name the same propositional variable with opposite
signs.
\end{remark*}

\begin{remark*}[Interpretation]
Complementary literals are the local contradiction pattern inside clauses and
terms. The pair \(P,\neg P\) is tied to one variable, not merely to two formulas
with opposite truth values under a particular assignment.
\end{remark*}

\begin{remark*}[Examples]
The literals \(P\) and \(\neg P\) are complementary, as are \(Q\) and
\(\neg Q\).
\end{remark*}

\begin{remark*}[Non-Examples]
The literals \(P\) and \(\neg Q\) are not complementary unless \(P=Q\).
\end{remark*}

\begin{dependencies}
\begin{itemize}
  \item \hyperref[def:literal-propositional-logic]{Literal}
\end{itemize}
\end{dependencies}
% END GENERATED ARTIFACT: def:complementary-literals-propositional-logic

% BEGIN GENERATED ARTIFACT: def:clause-propositional-logic
\begin{tcolorbox}[colback=propbox, colframe=propborder, arc=2pt,
  left=6pt, right=6pt, top=4pt, bottom=4pt,
  title={\small\textbf{Definition (Clause)}},
  fonttitle=\small\bfseries]
\begin{definition}[Clause]
\label{def:clause-propositional-logic}
A \emph{clause} is a finite disjunction of literals.
\end{definition}
\end{tcolorbox}

\begin{remark*}[Standard quantified statement]
A formula \(C\in\WFF\) is a clause if there are literals
\(\ell_1,\ldots,\ell_n\) with \(n\ge 1\) such that
\[
C=\ell_1\lor\cdots\lor\ell_n.
\]
The empty clause is treated separately as the empty disjunction.
\end{remark*}

\begin{remark*}[Definition predicate reading]
\[
\operatorname{Clause}(C)
\]
means that \(C\) is built by finitely disjoining literals.
\end{remark*}

\begin{remark*}[Interpretation]
A clause records a finite list of literal-level alternatives. It is satisfied
when at least one of its literals is true.
\end{remark*}

\begin{remark*}[Examples]
The formulas \(P\lor Q\), \(P\lor\neg Q\), and
\(\neg P\lor Q\lor R\) are clauses.
\end{remark*}

\begin{remark*}[Non-Examples]
The formula \(P\land Q\) is not a clause, because its top-level connective is a
conjunction. The formula \(\neg(P\lor Q)\) is not a clause, because its negation
has not been pushed down to the literal level.
\end{remark*}

\begin{dependencies}
\begin{itemize}
  \item \hyperref[def:literal-propositional-logic]{Literal}
\end{itemize}
\end{dependencies}
% END GENERATED ARTIFACT: def:clause-propositional-logic

% BEGIN GENERATED ARTIFACT: def:term-propositional-logic
\begin{tcolorbox}[colback=propbox, colframe=propborder, arc=2pt,
  left=6pt, right=6pt, top=4pt, bottom=4pt,
  title={\small\textbf{Definition (Term)}},
  fonttitle=\small\bfseries]
\begin{definition}[Term]
\label{def:term-propositional-logic}
A \emph{term}, also called a \emph{conjunction term}, is a finite conjunction
of literals.
\end{definition}
\end{tcolorbox}

\begin{remark*}[Standard quantified statement]
A formula \(T\in\WFF\) is a term if there are literals
\(\ell_1,\ldots,\ell_n\) with \(n\ge 1\) such that
\[
T=\ell_1\land\cdots\land\ell_n.
\]
The empty conjunction is treated separately.
\end{remark*}

\begin{remark*}[Definition predicate reading]
\[
\operatorname{Term}(T)
\]
means that \(T\) is built by finitely conjoining literals.
\end{remark*}

\begin{remark*}[Interpretation]
A term records a finite list of literal-level requirements. It is satisfied
when all of its literals are true.
\end{remark*}

\begin{remark*}[Examples]
The formulas \(P\land Q\), \(P\land\neg Q\), and
\(\neg P\land Q\land R\) are terms.
\end{remark*}

\begin{remark*}[Non-Examples]
The formula \(P\lor Q\) is not a term, because its top-level connective is a
disjunction. The formula \(\neg(P\land Q)\) is not a term, because its negation
has not been pushed down to the literal level.
\end{remark*}

\begin{dependencies}
\begin{itemize}
  \item \hyperref[def:literal-propositional-logic]{Literal}
\end{itemize}
\end{dependencies}
% END GENERATED ARTIFACT: def:term-propositional-logic

% BEGIN GENERATED ARTIFACT: def:empty-clause-propositional-logic
\begin{tcolorbox}[colback=propbox, colframe=propborder, arc=2pt,
  left=6pt, right=6pt, top=4pt, bottom=4pt,
  title={\small\textbf{Definition (Empty Clause)}},
  fonttitle=\small\bfseries]
\begin{definition}[Empty Clause]
\label{def:empty-clause-propositional-logic}
The \emph{empty clause} is the empty disjunction of literals. Semantically, it
is read as false.
\end{definition}
\end{tcolorbox}

\begin{remark*}[Standard quantified statement]
The empty clause is the disjunction over no literals and has truth value
\(\False\) under every truth assignment.
\end{remark*}

\begin{remark*}[Definition predicate reading]
\[
\operatorname{EmptyClause}(C)
\]
means that \(C\) is the clause with no disjuncts.
\end{remark*}

\begin{remark*}[Interpretation]
A nonempty clause is true when at least one listed literal is true. With no
literals listed, there is no witness to make the disjunction true, so the empty
clause is the false boundary case for clauses.
\end{remark*}

\begin{dependencies}
\begin{itemize}
  \item \hyperref[def:clause-propositional-logic]{Clause}
  \item \hyperref[def:contradiction-propositional-logic]{Contradiction}
\end{itemize}
\end{dependencies}
% END GENERATED ARTIFACT: def:empty-clause-propositional-logic

% BEGIN GENERATED ARTIFACT: def:empty-conjunction-propositional-logic
\begin{tcolorbox}[colback=propbox, colframe=propborder, arc=2pt,
  left=6pt, right=6pt, top=4pt, bottom=4pt,
  title={\small\textbf{Definition (Empty Conjunction)}},
  fonttitle=\small\bfseries]
\begin{definition}[Empty Conjunction]
\label{def:empty-conjunction-propositional-logic}
The \emph{empty conjunction} is the conjunction over no literals. Semantically,
it is read as true.
\end{definition}
\end{tcolorbox}

\begin{remark*}[Standard quantified statement]
The empty conjunction is the conjunction over no literals and has truth value
\(\True\) under every truth assignment.
\end{remark*}

\begin{remark*}[Definition predicate reading]
\[
\operatorname{EmptyConjunction}(T)
\]
means that \(T\) is the term with no conjuncts.
\end{remark*}

\begin{remark*}[Interpretation]
A nonempty term is true when every listed literal is true. With no literals
listed, there is no requirement to violate, so the empty conjunction is the true
boundary case for terms.
\end{remark*}

\begin{dependencies}
\begin{itemize}
  \item \hyperref[def:term-propositional-logic]{Term}
  \item \hyperref[def:tautology-propositional-logic]{Tautology}
\end{itemize}
\end{dependencies}
% END GENERATED ARTIFACT: def:empty-conjunction-propositional-logic
